% ===== §29 =====
\section{The Closed Loop Back into Intimate Relation}
\label{p24:sec:feedback}

\noindent
The capture described in the previous section is not a process that occurs elsewhere and leaves intimate relation untouched. It closes a loop: the value drawn from culture returns, as a force, into the intimate relation from which, at the last stage, it was drawn, and it alters that relation, eroding the very capacity for endogenous cultural generation that the second part described. This section traces the closing of that loop, the feedback by which the capture of intimate culture reaches back into the couple and changes what they can generate and who they can become.

\subsection*{Mediation: the supplying of ready-made forms}

The first moment of the feedback is \emph{mediation}. Increasingly, the forms through which intimate relation is conducted---its idioms, its rites, its expectations, its very sense of what love is and how it is expressed---are supplied ready-made by the culture industry and its successors, rather than generated by the couple from within. What romance is, how a proposal is made, how devotion is shown, what an anniversary requires, what a shared life should look like: these are furnished as scripts, and the couple is invited, at every turn, to adopt the supplied form in place of generating its own. The negotiation with the outer symbolic order that the phenomenology described has, under this feedback, tilted decisively toward substitution: the couple ceases to generate its own culture and instead consumes the culture supplied to it, so that the co-signification that should have been theirs is performed in advance, by the market that wrote the script. This is the colonization of intimate culture from within its own generation, and it is, in the terms of the previous part, the second misrecognition---the treatment of the bond's meaning as an exhibitable, convertible asset---installed as the ordinary condition of intimate life. This connects the diagnosis to the analysis of the symbolic order's colonization of the relational interior that the series has developed elsewhere, under the heading of the seepage of external regimes into the dyadic interior.

\subsection*{The deepening of the hollow subject}

The second moment is the deepening of a hollowing the series has diagnosed from its beginning. When the meaning of intimate relation is outsourced to consumer symbols---when the couple no longer generates its own forms but consumes the forms supplied---the relation loses its capacity to generate its own meaning, and with it the subject loses one of the last sites in which it might have come to be through a generation genuinely its own. The mother paper of this series diagnosed the modern condition as a crisis of symbol and subject, an inflation of signs accompanied by a hollowing of the subject, and the capture of intimate culture is the reaching of that crisis into the last interior \parencite{wan2024grb}. If the subject is generated in relation, as the whole of this paper has argued, and if the relation's capacity to generate is colonized by a market that supplies its meanings ready-made, then the subject is hollowed at the very site where it might have been filled, generated as the fitted consumer of supplied meanings rather than as the relational subject of an endogenous generation. The hollow subject is not only a consumer of public culture; it is now a consumer of its own intimate culture, deprived of the generation through which it might have become someone in particular, with someone in particular.

\subsection*{The artificial intelligence turn}

This deepening is redoubled by a development the mother paper marked and the present moment makes acute: the arrival of systems that can supply not only the forms of intimate culture but the generative act itself. Where the culture industry supplied the scripts and left the couple to perform them, there now arise systems that offer to perform the generation---to compose the message, to craft the endearment, to supply the words of devotion, to generate the very idiom that was to have been the couple's own. The final interiority that even the culture industry left to the couple---the act of generating their own forms, however much the raw materials were supplied---can now be outsourced to a machine. This is the capture reaching its innermost point: not merely the supplying of the forms the couple adopts, but the supplying of the generative act by which forms are made, so that even the coining of a private word, the making of a private meaning, is delegated to a system that owns the process and returns the product. The danger is precise, and it is not that such systems produce nothing of worth; it is that they offer to perform the one act---the endogenous generation of a shared meaning---on which, this paper has argued, the whole worth of intimate culture depends, and that a couple which delegates that act retains the form and loses exactly what made the form a culture rather than a purchase.

\subsection*{Alienation completed}

Read through the Hegelian line, the closed loop is the completion of alienation, and naming it as such gathers the section. Production, the first line established, is objectification: the going-out of the subject into a made world, which is the necessary path of its self-knowledge. Alienation is what objectification becomes when the made world is taken from its maker and turned against them---when what the subject produced stands over against it as an alien power, owned by another, ruling the very subject whose activity produced it. The closed loop is exactly this. The couple generates intimate culture; the culture is captured, its forms supplied and owned by the market, its generation colonized, at the limit its generative act itself delegated to a machine; and the captured culture returns to the couple as an alien power, a set of forms they consume rather than make, ruling their intimate life from a position they no longer occupy. What was to have been the couple's exteriorization---their going-out into a shared world in which they would know themselves and return to themselves enriched---becomes, under capture, an alienation: the shared world is owned by another, and it returns to them not as their own work reflecting them back to themselves but as a foreign power shaping them to its ends. The exteriorized culture that should have been the couple's self-knowledge becomes the instrument of their subjection. This is the darkest point of the paper, and it is the point at which the argument must turn: for if the framework of Part One is right, the alienation cannot be total, and the turn from diagnosis to the conditions of recovery begins with the principle the next section states---that the material conditions constrain but do not determine, and that the capacity for endogenous generation, though colonized, is never wholly destroyed.
